\sisection{S6. Experimental resources, endpoints and boundaries}

This section defines the external endpoints and records the sensitivity checks and negative
results behind the main-text comparisons. Values come from frozen files under \path{results/},
and each table names its source. These analyses are exploratory and post hoc as recorded in their
summary files. ``Frozen'' identifies a reproducible analysis definition; it does not imply
prospective registration. The homology controls in S6.3 were added after the first comparison.
Except for the explicitly corrected PDSP family, nominal probabilities are descriptive and are
not adjusted across the wider search.

\sisubsection{S6.1 The frozen kinase endpoint}

The endpoint uses no docking score. DAVIS, PKIS2 and PKIS1 share 20 kinase labels, giving 190
unordered target pairs. Within each panel we two-way centred the raw activity matrix, correlated
its target columns, and called a pair positive when its centred correlation fell in the upper
decile of that panel. A pair is positive for the endpoint when it is positive in at least two of
the three panels, which labels 15 of the 190 pairs; 7 of those 15 are positive in all three
panels. Panel sizes are 72 compounds for DAVIS, 645 for PKIS2 and 360 for PKIS1
(\path{results/replicated_pair_retrieval/summary.json}).

The docking predictor is built on a reference support that excludes the assay chemistry. Any
DOCKSTRING row whose 14-character Standard InChI connectivity block appears in one of the three
panels was removed before any correlation was computed: 1{,}050 blocks and 481 rows, leaving
259{,}579 of 260{,}060 rows. Positive scores are clipped to zero; the unclipped sensitivity is in
S6.2.

Inference permutes the 20 docking target labels and holds the experimental endpoint and every
structural control fixed. Each draw is one complete relabelling applied to the rows and columns
of the docking map, so all 190 edges move together. Probabilities are one-sided for a positive
association with the add-one correction, over 50{,}000 permutations at seed 20260803. Because
the analysis is post hoc, they are descriptive directional tail areas. With
50{,}000 draws the smallest attainable value is $2.0\times10^{-5}$.

\begin{table}[htbp]
\centering
\small
\setlength{\tabcolsep}{5pt}
\caption[Predictors of the frozen 190-pair kinase endpoint]{Predictors of the frozen 190-pair kinase endpoint, 15 positives, evaluated on identical
pairs. Probabilities are one-sided target-label permutation values over 50{,}000 draws at seed
20260803. Source: \path{results/klifs_pocket_control/locked_endpoint_retrieval.csv}.}
\label{tab:s6-frozen-predictors}
\begin{tabular}{lrrrr}
\toprule
Predictor & ROC-AUC & $p$ & Average precision & $p$ \\
\midrule
Raw Vina target map            & 0.5966 & 0.155  & 0.2355 & 0.0113 \\
Centred Vina target map        & 0.7105 & 0.0069 & 0.3015 & 0.0014 \\
Receptor-domain sequence identity & 0.7181 & 0.0070 & 0.4263 & $2.0\times10^{-5}$ \\
KLIFS 85-residue pocket identity  & 0.7082 & 0.0103 & 0.4864 & $2.0\times10^{-5}$ \\
Same KLIFS family              & 0.5971 & 0.0012 & 0.1916 & 0.0021 \\
Same KLIFS group               & 0.6514 & 0.0431 & 0.1265 & 0.0160 \\
\bottomrule
\end{tabular}
\end{table}

The paired increments from raw to centred docking, computed under the same relabellings, are
$+0.1139$ in ROC-AUC and $+0.0660$ in average precision, with target-label probabilities 0.0334
and 0.0166 and Holm-adjusted values 0.0334 and 0.0333 across the two metrics. The 95\% null
intervals for those paired differences are $-0.108$ to $+0.123$ and $-0.054$ to $+0.058$.
These adjusted values describe the two selected-endpoint directional contrasts; they are not
confirmatory evidence across the threshold search.

Two readings of Table~\ref{tab:s6-frozen-predictors} matter. Centred docking reaches sequence
identity on ROC-AUC, 0.7105 against 0.7181, and stays below it on average precision, 0.3015
against 0.4263; pocket identity is the strongest predictor on average precision at 0.4864.
Centred docking does not pass either structural baseline. The binary annotations behave
differently from the continuous ones under this null: same-family membership has ROC-AUC 0.5971,
close to chance, yet $p=0.0012$, because only 4 of 190 pairs are same-family and the
relabelling distribution of that statistic is coarse. A small probability on a two-valued
predictor here reports rarity of the configuration, not separation.

Centred Vina geometry is not itself a stand-in for the annotations it is being compared with. Its
rank association with pocket identity is 0.0669 ($p=0.175$), with same-family 0.1123 ($p=0.0637$)
and with same-group 0.0307 ($p=0.266$)
(\path{results/klifs_pocket_control/annotation_associations.csv}).

\sisubsection{S6.2 Threshold and transform sensitivity of the kinase gain}

The upper decile is one cut of a continuous experimental map. Table~\ref{tab:s6-threshold} sweeps
the cut from 5\% to 30\%. The declared omnibus covers the five thresholds 0.05, 0.10, 0.15, 0.20
and 0.25: the mean paired gain is $+0.0723$ in ROC-AUC at $p=0.0605$ and $+0.0323$ in average
precision at $p=0.0759$, over 50{,}000 draws at seed 20360803, with 95\% null intervals $-0.080$
to $+0.093$ and $-0.049$ to $+0.051$. Neither omnibus reaches 0.05. The gain is an upper-tail
result and is not stable across the whole threshold family: at 0.25 the ROC-AUC gain is $+0.0003$
and the average-precision gain is $-0.0125$.

\begin{table}[htbp]
\centering
\small
\setlength{\tabcolsep}{4pt}
\caption[Threshold family for the experimental upper tail]{Threshold family for the experimental upper tail. The 10\% row is the frozen endpoint.
Rows at 0.075, 0.125 and 0.30 are outside the declared omnibus family. Source:
\path{results/replicated_pair_retrieval/threshold_sensitivity.csv}.}
\label{tab:s6-threshold}
\begin{tabular}{rrrrrrrr}
\toprule
Upper-tail & Positive & \multicolumn{3}{c}{ROC-AUC} & \multicolumn{3}{c}{Average precision} \\
\cmidrule(lr){3-5}\cmidrule(lr){6-8}
fraction & pairs & raw & centred & $\Delta$ & raw & centred & $\Delta$ \\
\midrule
0.050 &  9 & 0.6262 & 0.7041 & $+0.0780$ & 0.1922 & 0.2090 & $+0.0168$ \\
0.075 & 14 & 0.5913 & 0.6952 & $+0.1039$ & 0.2398 & 0.2967 & $+0.0569$ \\
0.100 & 15 & 0.5966 & 0.7105 & $+0.1139$ & 0.2355 & 0.3015 & $+0.0660$ \\
0.125 & 21 & 0.6475 & 0.7588 & $+0.1113$ & 0.2580 & 0.3431 & $+0.0850$ \\
0.150 & 24 & 0.6606 & 0.7636 & $+0.1029$ & 0.3109 & 0.3537 & $+0.0427$ \\
0.200 & 35 & 0.6431 & 0.7095 & $+0.0664$ & 0.3290 & 0.3777 & $+0.0487$ \\
0.250 & 47 & 0.6535 & 0.6538 & $+0.0003$ & 0.3992 & 0.3868 & $-0.0125$ \\
0.300 & 56 & 0.6499 & 0.6530 & $+0.0031$ & 0.4353 & 0.4329 & $-0.0024$ \\
\bottomrule
\end{tabular}
\end{table}

The endpoint also depends on how the experimental matrix is preprocessed before its target
columns are correlated. Table~\ref{tab:s6-transform} gives three preprocessing choices. Centring
then correlating is the primary rule and reproduces the primary labels exactly. Column
$z$-scoring before centring keeps 15 positives but changes which pairs they are, with Jaccard
0.429 against the primary label set, and it raises the paired gain to $+0.2510$ in ROC-AUC.
Column inverse-normal ranks before centring give 18 positives and Jaccard 0.500. The paired gain
is positive under all three rules, and its size moves by a factor of 2.2 in ROC-AUC. Clipping is
not what carries the result: with positive scores left unclipped, the endpoint gives raw 0.5943
and centred 0.7101 in ROC-AUC, and raw 0.2335 and centred 0.3006 in average precision.

\begin{table}[htbp]
\centering
\small
\setlength{\tabcolsep}{4pt}
\caption[Experimental-transform sensitivity of the frozen endpoint]{Experimental-transform sensitivity of the frozen endpoint. Jaccard is the overlap of the
positive-pair set with the primary labels. Source:
\path{results/replicated_pair_retrieval/summary.json}, key
\texttt{experimental\_transform\_sensitivity}.}
\label{tab:s6-transform}
\begin{tabular}{lrrrrrrrr}
\toprule
Experimental & Pos. & \multicolumn{3}{c}{ROC-AUC} & \multicolumn{3}{c}{Average precision} & \\
\cmidrule(lr){3-5}\cmidrule(lr){6-8}
transform & pairs & raw & centred & $p(\Delta)$ & raw & centred & $p(\Delta)$ & Jaccard \\
\midrule
Centre, then correlate       & 15 & 0.5966 & 0.7105 & 0.0344   & 0.2355 & 0.3015 & 0.0171   & 1.000 \\
$z$-score, then centre       & 15 & 0.5170 & 0.7680 & 0.00034  & 0.2233 & 0.3994 & 0.00032  & 0.429 \\
Rank-normal, then centre     & 18 & 0.5346 & 0.7109 & 0.0018   & 0.1576 & 0.2708 & 0.0020   & 0.500 \\
\bottomrule
\end{tabular}
\end{table}

The first row reproduces the selected endpoint's observed statistics. Its tail areas differ
slightly from S6.1 because this transform-sensitivity run used an independent permutation stream
(seed 20460803 rather than 20260803); the interpretation is unchanged.

The remaining sensitivities hold the endpoint fixed and vary the support. Deleting one target at
a time leaves the ROC-AUC gain positive in all 20 deletions, from $+0.0294$ (KDR omitted) to
$+0.1659$ (AKT1 omitted), and the average-precision gain positive in all 20, from $+0.0049$ to
$+0.1003$ (\path{results/replicated_pair_retrieval/target_jackknife.csv}). Five independent
15{,}000-ligand reference supports at seeds 11, 29, 47, 71 and 97 give ROC-AUC gains from
$+0.1105$ to $+0.1185$ and average-precision gains from $+0.0326$ to $+0.0792$. Rebuilding the
endpoint from two panels and testing on the third gives ROC-AUC gains from $+0.0508$ to $+0.0960$
and average-precision gains from $-0.0128$ to $+0.0677$, so one of the three held-out
constructions loses on average precision
(\path{results/replicated_pair_retrieval/support_sensitivities.csv}).

\sisubsection{S6.3 Homology-restricted permutation}

This is the weakest structural point in the paper and we state it in full. The permutation of
S6.1 lets any docking target label move to any position. Kinases in one KLIFS group are more
alike than kinases from different groups, so an unrestricted null can be broken by group
structure alone. The restricted scheme allows a target label to exchange only with targets in the
same KLIFS group. Every draw is still a complete bijection of all 20 labels, and the experimental
endpoint, sequence identity and pocket identity stay fixed.

\begin{table}[htbp]
\centering
\small
\caption[Exchangeability blocks of the fixed 20-kinase panel]{Exchangeability blocks of the fixed 20-kinase panel. Source:
\path{results/strict_klifs_group_qap/exchangeability_blocks.csv}.}
\label{tab:s6-blocks}
\begin{tabular}{llrrr}
\toprule
KLIFS group & Targets & Size & Movable & Label orders \\
\midrule
TK    & ABL1, CSF1R, EGFR, FGFR1, IGF1R, JAK2, & 11 & yes & 39{,}916{,}800 \\
      & KDR, KIT, LCK, MET, SRC & & & \\
AGC   & AKT1, AKT2, ROCK1 & 3 & yes & 6 \\
CMGC  & CDK2, MAPK1, MAPK14 & 3 & yes & 6 \\
STE   & MAP2K1 & 1 & no & 1 \\
CAMK  & MAPKAPK2 & 1 & no & 1 \\
Other & PLK1 & 1 & no & 1 \\
\bottomrule
\end{tabular}
\end{table}

The restricted space has $11!\times 3!\times 3! = 1{,}437{,}004{,}800$ orders. Of the 20 target
labels, 17 can move and 3 cannot. Of the 190 pair positions, 187 can receive a different
predictor value and 3 cannot, those being the pairs among the singleton groups. TK dominates
what is left: 154 of 190 positions (81.1\%) involve a TK target, so the test is mostly a test of
the TK block. The primary null drew 100{,}000 group-restricted orders at seed 20260805, of which
99{,}997 were unique.

\begin{table}[htbp]
\centering
\small
\setlength{\tabcolsep}{4pt}
\caption[Within-KLIFS-group restricted permutation]{Within-KLIFS-group restricted permutation of centred Vina target geometry. Probabilities
are one-sided for a positive association with the add-one correction. Source:
\path{results/strict_klifs_group_qap/restricted_partial_qap.csv}.}
\label{tab:s6-restricted}
\begin{tabular}{llrrrr}
\toprule
Experimental map & Control & $\rho$ & $p$ & Null median & Null 95\% interval \\
\midrule
Three-panel mean & none & 0.2676 & 0.167 & 0.1346 & $-0.019$ to $0.434$ \\
Three-panel mean & sequence $+$ pocket & 0.2247 & 0.197 & 0.1030 & $-0.063$ to $0.415$ \\
KiRHub & none & 0.3216 & 0.117 & 0.1867 & $0.035$ to $0.422$ \\
KiRHub & sequence $+$ pocket & 0.2799 & 0.166 & 0.1614 & $-0.019$ to $0.419$ \\
\bottomrule
\end{tabular}
\end{table}

The two sequence-and-pocket-adjusted rows have exact counterparts under the unrestricted null of
S6.1, at the same point estimates: $p=0.0140$ for the three-panel mean and $p=0.0068$ for KiRHub
(\path{results/klifs_pocket_control/partial_geometry_qap.csv}). Restricting relabelling to KLIFS
groups moves both to 0.197 and 0.166.

The restricted null is broad enough that every observed value falls below its one-sided 5\%
critical value. Relabelling only the TK block gives much the same result, while exact enumeration
of the two three-target blocks has only 36 orders and little resolution. A stricter within-family
scheme is not useful here: only four pairs can exchange, leaving 16 orders and a smallest possible
one-sided probability of 0.0625. Full diagnostics are in
\path{results/strict_klifs_group_qap/summary.json}.

Homology-preserving nulls were also run on the two secondary-pharmacology resources. PDSP was a
post-hoc analysis. Under within-family relabelling, its residual-minus-raw contrasts had
two-sided probabilities of 0.0134 for continuous agreement, 0.2155 for top-decile ROC-AUC and
0.0209 for best-partner recall. Their Holm-adjusted values across the full 12-test PDSP family
were 0.1016, 0.2560 and 0.1254; the corresponding Bonferroni values were 0.1608, 1.000 and
0.2508. For the residual increment over the sequence-plus-family baseline, the smallest
within-family values were 0.0095 before adjustment, 0.0855 after Holm adjustment and 0.1140 after
Bonferroni adjustment. We therefore treat the PDSP homology pattern as exploratory. On the
Novartis panel, whose directional tests remain one-sided, the residual-minus-raw gain did not
persist on any of the three endpoints. The residual map itself remained associated with each
endpoint under the same null. Exact contrasts are in
\path{results/pdsp_counterscreen_retrieval/paired_qap.csv} and
\path{results/spd_external_validation/summary.json}.

The PDSP point estimates beyond a sequence-plus-family baseline also depend on docking support,
as detailed in S6.4. On the Novartis panel, evidence beyond those controls was endpoint-specific:
it remained for the residual map at 10\,$\mu$M, but the raw map also retained partial agreement.
The signal therefore appears in the broader docking geometry and is not unique to centring.
Kinase-group membership remains a viable explanation of the kinase result. The restricted test
also assumes exchangeability within each KLIFS group; groups still differ in family, pocket,
construct and assay context.

\sisubsection{S6.4 PDSP support, censoring and the counterscreen decision}

The PDSP $K_\mathrm{i}$ download holds 98{,}678 rows. The strict contract keeps human rows from
the allowed recombinant or cloned source vocabulary, exact relations only, positive $K_\mathrm{i}$
and a valid structure, then takes the median $\mathrm{p}K_\mathrm{i}$ per Standard InChIKey and
target: 3{,}722 source rows collapse to 1{,}560 compound-target cells over 481 compounds and 27
exported targets. An experimental edge requires at least 10 shared compounds, which gives the
primary graph of 148 target pairs over 21 targets. The Certified Data subset is smaller and
better documented: 44 compounds and 317 cells over 20 exported targets, whose graph has 78 pairs
over 15 targets.
The retained values come from heterogeneous publications and assay conditions. Taking medians
within compound--target cells and using rank correlations do not remove study or batch effects.

We removed PDSP chemistry from the docking reference before building any predictor. Compounds in
the exact-only endpoint share 63 full Standard InChIKeys and 116 connectivity blocks with
Docking-44. The censor-aware analyses use additional structures, so the conservative mask was
built from every PDSP structure entering any reported analysis. It removes 128 Docking-44 rows
and leaves 12{,}523. The primary geometry uses this de-overlapped support and fills the remaining
missing docking scores by target-column means. The exclusion list is
\path{results/pdsp_counterscreen_retrieval/excluded_docking44_ligand_ids.csv}.

\begin{table}[htbp]
\centering
\small
\setlength{\tabcolsep}{5pt}
\caption[Raw and residual target maps against the PDSP map]{Raw and residual Docking-44 target maps against the PDSP experimental target-pair map.
Top-decile labels are the upper 10\% of experimental edges. Sources:
\path{results/pdsp_counterscreen_retrieval/global_retrieval_metrics.csv} and
\path{results/pdsp_counterscreen_retrieval/certified_global_metrics.csv}.}
\label{tab:s6-pdsp}
\begin{tabular}{llrrrrr}
\toprule
Support & Map & Pairs & Positives & Spearman & ROC-AUC & Average precision \\
\midrule
Primary   & raw      & 148 & 15 & 0.0439 & 0.4747 & 0.1501 \\
Primary   & residual & 148 & 15 & 0.2953 & 0.6812 & 0.1771 \\
Certified & raw      &  78 &  8 & 0.0744 & 0.3786 & 0.0921 \\
Certified & residual &  78 &  8 & 0.4248 & 0.7607 & 0.2330 \\
\bottomrule
\end{tabular}
\end{table}

The PDSP analysis was post hoc, so its paired QAP uses two-sided probabilities. The correction
family contains two contrasts, three metrics and two relabelling schemes, for 12 tests. Under
unrestricted relabelling, the residual-minus-raw contrast has two-sided values of 0.0020 for
continuous agreement, 0.0048 for top-decile ROC-AUC and 0.0027 for best-partner recall. The
corresponding Holm values are 0.0240, 0.0480 and 0.0297; the Bonferroni values are 0.0240, 0.0576
and 0.0324. Adding the residual map to the sequence-plus-family baseline raises continuous
agreement by 0.0963 on this support, but its two-sided value of 0.0127 becomes 0.1016 by Holm and
0.1524 by Bonferroni. The within-family results are summarised in S6.3. These QAPs relabel an
already estimated edge map: they preserve dependence among target pairs but do not propagate
uncertainty from estimating each experimental edge on a small, pair-specific compound set.

The docking support is consequential. After the same overlap exclusion, only 9{,}221 Docking-44
rows have all 44 scores observed. This complete-case subset is chemically shifted relative to
the full support, so it is not a pure comparison of imputation rules. Of the 3{,}302 rows removed,
3{,}300 are missing the CHRM2 score column (PDB 4mqs). On it, raw and residual
continuous agreement are $-0.0987$ and 0.0463, and adding the residual to the
sequence-plus-family baseline changes agreement by $-0.0381$. For the continuous
residual-minus-raw contrast, the two-sided QAP value is 0.0713, with Holm 0.7843 and Bonferroni
0.8556 across its 12-test sensitivity family. Thus the favorable full-support estimate does not
transfer to this chemically different support. Complete results are in
\path{results/pdsp_counterscreen_retrieval/complete_case_support_metrics.csv} and
\path{results/pdsp_counterscreen_retrieval/complete_case_support_qap.csv}.

Censoring changes the experimental graph as well. The exact-only graph conditions on both targets
having a quantified $K_\mathrm{i}$ for the same compounds. Restoring censored observations at
their reported bounds expands the primary graph from 148 to 191 pairs and moves residual
agreement from 0.2953 to 0.0759. Binary recoding of quantified versus censored cells gives 0.0512.
The Certified graph shows the same weakening. Table~\ref{tab:s6-pdsp-censor} gives both docking
maps under the restored-censor endpoints.

\begin{table}[htbp]
\centering
\small
\setlength{\tabcolsep}{5pt}
\caption[PDSP censor-aware sensitivity]{PDSP censor-aware sensitivity. The final column is
top-decile ROC-AUC for the residual map. Source:
\path{results/pdsp_counterscreen_retrieval/censor_sensitivity_metrics.csv}.}
\label{tab:s6-pdsp-censor}
\begin{tabular}{lrrrr}
\toprule
Support and endpoint & Pairs & Raw $\rho$ & Residual $\rho$ & Residual ROC-AUC \\
\midrule
All allowed, reported bounds & 191 & $-0.1536$ & 0.0759 & 0.6140 \\
All allowed, quantified/censored & 191 & $-0.1643$ & 0.0512 & 0.6544 \\
Certified, reported bounds & 171 & 0.0041 & 0.1499 & 0.7222 \\
Certified, quantified/censored & 171 & $-0.0251$ & 0.0568 & 0.6107 \\
\bottomrule
\end{tabular}
\end{table}

Best-partner retrieval also changes with the censoring rule. With reported-bound insertion, the
residual map recalls 0.25 versus 0.20 for the raw map. After binary recoding, the residual map
recalls 0.20 versus 0.35 for the raw map. On the exact-only graph the values are 0.611 and 0.389.
The sequence-plus-family baseline is 0.556 there; adding the residual map to that baseline reaches
0.667. These are distinct predictors, and the fusion value is not the residual-only value. Exact
targetwise results are in
\path{results/pdsp_counterscreen_retrieval/summary.json} and the companion CSV files.

The minimum compound support for an experimental edge was also varied after seeing the primary
result. Table~\ref{tab:s6-pdsp-support} is therefore a descriptive sweep, not a separate
confirmatory test. Changing the cutoff changes the edge set and, at the higher cutoffs, the target
set. Residual agreement remains above raw agreement in every row, but its size varies, as does the
increment over the sequence-plus-family baseline.

\begin{table}[htbp]
\centering
\small
\setlength{\tabcolsep}{6pt}
\caption[Minimum pair-support sensitivity on PDSP]{Minimum shared compounds per PDSP edge.
The final column is the continuous-agreement increment from adding the residual map to an
equal-rank sequence-plus-family baseline. Source:
\path{results/pdsp_counterscreen_retrieval/minimum_pair_support_sensitivity.csv}.}
\label{tab:s6-pdsp-support}
\begin{tabular}{rrrrr}
\toprule
Minimum compounds & Pairs & Raw $\rho$ & Residual $\rho$ & Fusion increment \\
\midrule
10 & 148 &  0.0439 & 0.2953 & 0.0963 \\
12 & 140 &  0.0266 & 0.2844 & 0.0831 \\
15 & 128 &  0.0352 & 0.2560 & 0.0670 \\
20 &  98 &  0.0053 & 0.2017 & 0.0429 \\
25 &  72 &  0.0993 & 0.4327 & 0.1605 \\
30 &  51 & $-0.0410$ & 0.2253 & 0.0807 \\
\bottomrule
\end{tabular}
\end{table}

A support control isolates the effect of pooling different compounds into different edges. There
are 113 five-target blocks in which at least 20 compounds are measured on all five targets, with
a median of 21 complete compounds. Across those blocks the median residual-minus-raw difference
in top-20\% ROC-AUC is $+0.0625$, and 61.1\% of blocks are positive.

The frozen go/no-go rule required a residual-versus-raw gain under aligned-target permutation,
most target deletions, the censor-aware endpoints, the Certified-only support and a majority of
identical-support blocks. The recorded verdict is
\texttt{NO\_GO\_GENERAL\_COUNTERSCREEN}. The remaining pattern is narrow: the residual map is
more similar to the exact-only PDSP graph when both targets have a quantifiable $K_\mathrm{i}$
for shared compounds. It is not a general ordering for deciding which
counterscreens to run, and we do not present it as one.

\sisubsection{S6.5 SPD endpoints, censoring and normalisation comparison}

The Novartis Safety Profiling Database export covers 12 human targets. For each target, the
analysis uses the candidate direct binding or inhibition assay group with the largest coverage;
ties are broken by the smaller group identifier.
ESR1, PGR and PTGS2 had a second campaign available
(\path{results/spd_external_validation/assay_selection.csv}). The panel holds 1{,}975 full
Standard InChIKeys across 1{,}925 connectivity blocks, and none of the 1{,}975 keys matches a
DOCKSTRING compound exactly. We removed all 932 shared connectivity blocks from the docking side
anyway, dropping 1{,}364 DOCKSTRING rows. Of 11{,}126 selected activity rows, 10{,}188 are
right-censored, which is 91.6\%; 938 rows carry an exact value. An experimental edge requires 40
shared compounds, giving 59 target pairs.

\begin{table}[htbp]
\centering
\small
\setlength{\tabcolsep}{5pt}
\caption[Docking and SPD target-pair map agreement]{Agreement between DOCKSTRING target maps and SPD experimental target-pair maps, on 59
pairs, under both experimental transforms. The docking side is the 12-target local map built on
the non-overlapping DOCKSTRING support. Source:
\path{results/spd_external_validation/geometry_metrics.csv}.}
\label{tab:s6-spd}
\begin{tabular}{llrrr}
\toprule
Experimental endpoint & Assay map & Raw $\rho$ & Residual $\rho$ & Difference \\
\midrule
Reported bound ranks     & centred & 0.1488 & 0.3723 & $+0.2235$ \\
Reported bound ranks     & raw     & 0.3826 & 0.5139 & $+0.1313$ \\
Active at 10\,$\mu$M     & centred & 0.3479 & 0.3884 & $+0.0404$ \\
Active at 10\,$\mu$M     & raw     & 0.2786 & 0.5881 & $+0.3095$ \\
Active at 30\,$\mu$M     & centred & 0.1804 & 0.4513 & $+0.2709$ \\
Active at 30\,$\mu$M     & raw     & 0.1635 & 0.3344 & $+0.1709$ \\
\bottomrule
\end{tabular}
\end{table}

Against the centred assay map, the residual representation is associated with all three endpoints
under unrestricted target relabelling ($p=0.0124$, $0.0052$ and $0.0058$). Its gain over raw is
supported for reported-bound ranks and 30\,$\mu$M ($p=0.0376$ and $0.0206$), but not for
10\,$\mu$M ($p=0.3811$). That endpoint is a counterexample within the same resource. After rank
control for sequence, family, support and coverage, residual partial agreement on reported-bound
ranks is 0.2405 at $p=0.0572$; we therefore make no beyond-control claim for this endpoint.
Family-preserving results are in S6.3.

The observation mask is target-specific, so we also built the largest outcome-blind complete
rectangle, which fixes the compound set across targets: 102 compounds measured on the same 11
targets, all except EGFR, giving 55 pairs. There raw agreement is 0.0552 and residual agreement
is 0.5752, a difference of $+0.5200$, with raw partial rank $-0.2201$ and residual partial rank
$+0.3706$. Deleting one target at a time from the primary endpoint leaves all 12 differences
positive, from $+0.1610$ (DRD2 omitted) to $+0.3107$ (ADORA2A omitted)
(\path{results/spd_external_validation/target_jackknife.csv}). The exact-only surface is too thin
to carry anything: at minimum support 5 it has 28 pairs with raw 0.3905 and residual 0.4037,
and at minimum support 10 it has 16 pairs with raw $-0.2208$ and residual $-0.0059$.

Two-way centring is not the only within-ligand normalisation that works, and it is not the best
one on this panel. Table~\ref{tab:s6-norm} compares five representations on the same three
endpoints and the same de-leaked support.

\begin{table}[htbp]
\centering
\small
\setlength{\tabcolsep}{5pt}
\caption[Normalisation controls on the SPD panel]{Normalisation controls on the SPD panel, 59 pairs per cell, Spearman agreement with the
fixed centred assay map. Source: \path{results/spd_normalization_controls/metrics.csv}.}
\label{tab:s6-norm}
\begin{tabular}{lrrr}
\toprule
Docking representation & Bound ranks & 10\,$\mu$M & 30\,$\mu$M \\
\midrule
Raw, target-standardised          & 0.1488 & 0.3479 & 0.1804 \\
Standard centring                 & 0.3723 & 0.3884 & 0.4513 \\
Within-ligand ordinal             & 0.3663 & 0.2877 & 0.3633 \\
Ligand efficiency, raw            & 0.2849 & 0.2933 & 0.2642 \\
Ligand efficiency, centred        & 0.4494 & 0.2828 & 0.4445 \\
\bottomrule
\end{tabular}
\end{table}

Ligand-efficiency centring is best on reported-bound ranks, while standard centring is best at
10\,$\mu$M. No paired contrast against standard centring reaches $p=0.05$. The ligand-efficiency
advantage is stable across target deletions only for the reported-bound endpoint. Exact
permutation and deletion results are in
\path{results/spd_normalization_controls/target_label_qap.csv} and
\path{results/spd_normalization_controls/target_jackknife.csv}. The choice of normalisation is a
workflow decision that should be reported; its ranking on one panel does not transfer.

The ordinal transform is the clearest case of non-transfer. It is worse than the column-$z$ plus
row-centred comparator on all four kinase panels and in all 20 target deletions. The equal-weight
difference is $-0.0424$ (two-sided $p=0.0560$). Its only advantage over this comparator is on the
thresholded kinase endpoint, ROC-AUC 0.7303 versus 0.7181, and the paired permutation does not
support that difference ($p=0.321$).

\begin{table}[htbp]
\centering
\small
\caption[Within-ligand ordinal geometry on four kinase panels]{Within-ligand ordinal geometry against the column-$z$ plus row-centred comparator on four kinase panels, Spearman
agreement with each panel's fixed centred target map. Source:
\path{results/ordinal_docking_geometry/representation_metrics.csv} and
\path{results/ordinal_docking_geometry/paired_target_label_qap.csv}.}
\label{tab:s6-ordinal}
\begin{tabular}{lrrrr}
\toprule
Panel & Column-$z$ + row-centred & Ordinal & Difference & Two-sided $p$ \\
\midrule
DAVIS  & 0.2925 & 0.2526 & $-0.0399$ & 0.0909 \\
PKIS2  & 0.3003 & 0.2462 & $-0.0541$ & 0.0302 \\
PKIS1  & 0.1854 & 0.1419 & $-0.0435$ & 0.0712 \\
KiRHub & 0.3361 & 0.3040 & $-0.0321$ & 0.1386 \\
Equal-weight mean & 0.2786 & 0.2362 & $-0.0424$ & 0.0560 \\
\bottomrule
\end{tabular}
\end{table}

\sisubsection{S6.6 The KiRHub non-replication}

KiRHub supplies 92 drugs by 20 kinase targets, 1{,}840 cells with none missing, as residual
activity at a single 1\,$\mu$M drug concentration and $K_\mathrm{m}$ ATP. The measurement has a
hard ceiling: 372 cells (20.2\%) sit at exactly 100 and 9 cells (0.49\%) at exactly 0, and the
per-target ceiling fraction runs from 0.076 (FGFR1) to 0.511 (AKT2), with a median of 0.185
(\path{results/kirhub_external_validation/target_ceiling_audit.csv}).

The docking comparison uses a KiRHub-only endpoint: the upper decile of the centred KiRHub
target-pair map, 19 positives among 190 pairs. Raw docking gives ROC-AUC 0.7510 and average
precision 0.3312; centred docking gives 0.7436 and 0.2895. Both paired differences are negative,
and neither supports a positive gain. Exact probabilities and null distributions are in
\path{results/kirhub_external_validation/summary.json}. This is a non-replication of the
centred-minus-raw gain on a fourth kinase panel.

No choice of docking-side transform recovers it (Table~\ref{tab:s6-kirhub-transform}). Removing
target offsets alone leaves the correlation unchanged, as it must. Removing the per-ligand mean
loses 0.042 in average precision. Column $z$-scoring or inverse-normal ranks before centring
recovers ROC-AUC to 0.753 and 0.754 while leaving average precision at 0.282 and 0.285. Leaving
positive scores unclipped changes ROC-AUC by less than 0.003 in either representation.

\begin{table}[htbp]
\centering
\small
\caption[Docking-side transforms against the KiRHub endpoint]{Docking-side transforms against the KiRHub-only endpoint, 190 pairs with 19 positives.
Source: \path{results/kirhub_external_validation/transform_sensitivity.csv}.}
\label{tab:s6-kirhub-transform}
\begin{tabular}{lrr}
\toprule
Docking transform & ROC-AUC & Average precision \\
\midrule
Raw                                & 0.7510 & 0.3312 \\
Target centring only               & 0.7510 & 0.3312 \\
Per-ligand centring                & 0.7436 & 0.2895 \\
Two-way centring                   & 0.7436 & 0.2895 \\
$z$-score, then centre             & 0.7528 & 0.2821 \\
Rank-normal, then centre           & 0.7544 & 0.2845 \\
Raw, unclipped                     & 0.7485 & 0.3273 \\
Two-way centred, unclipped         & 0.7433 & 0.2892 \\
\bottomrule
\end{tabular}
\end{table}

Chemical independence from the earlier panels can only be audited by name. KiRHub tables and the
portal give no machine-readable structures for the 92 drugs, so no InChIKey, connectivity or
scaffold comparison is possible, and the DOCKSTRING reference support could not be de-leaked
against KiRHub compounds the way it was for DAVIS, PKIS2 and PKIS1. Exact normalised names match
the DAVIS 72-drug panel in 11 cases. Removing those 11 leaves 81 drugs, and the KiRHub
experimental geometry then agrees with the three-panel mean at 0.8166 rather than 0.8340, and
recovers the frozen endpoint at ROC-AUC 0.9699 rather than 0.9634 and average precision 0.7353
rather than 0.6937 (\path{results/kirhub_external_validation/drug_name_overlap_sensitivity.csv}).
The overlap does not carry the KiRHub experimental result, and the overlap sensitivity is a check
on the experimental side only; it is not a structure-level independence claim. Using
CDK2--cyclin E instead of CDK2--cyclin A moves the three-panel geometry agreement from 0.8340 to
0.8335.

One distinction must stay sharp. Centring the KiRHub \emph{assay} matrix helps: recovery of the
frozen endpoint rises from ROC-AUC 0.8034 to 0.9634 and average precision from 0.5971 to 0.6937.
Centring the \emph{docking} matrix against the KiRHub endpoint does not. These are different
sides of the comparison and the first does not support the second.

\sisubsection{S6.7 Ligand-level retrieval boundary}

The claims above concern target-pair maps. Ranking targets for one ligand is a different
estimand, and on that estimand centring is worse. The frozen benchmark takes ChEMBL records with
an exact relation for $K_\mathrm{i}$, $K_\mathrm{d}$, IC$_{50}$ or EC$_{50}$, medians them per
compound-target cell, and calls a cell active at $\mathrm{pChEMBL} \geq 6.0$. Against the
44-target Docking-44 panel this gives 500 active compound-target pairs over 190 compounds and 39
represented targets, spread over 166 Butina clusters. Compounds carry a median of 1 active target
and at most 14.

\begin{table}[htbp]
\centering
\small
\setlength{\tabcolsep}{5pt}
\caption[Single-ligand target retrieval on the frozen ChEMBL active set]{Single-ligand target retrieval on the frozen ChEMBL active set, compound-balanced
weighting. Lower mean rank percentile is better; higher top-$k$ is better. Intervals are
Butina-cluster bootstrap 95\% intervals over 5{,}000 replicates. Source:
\path{results/residual_mie_boundary_audit/observed_metrics.csv} and
\path{results/residual_mie_boundary_audit/paired_contrasts.csv}.}
\label{tab:s6-ligand}
\begin{tabular}{lrrrl}
\toprule
Metric & Absolute Vina & Two-way residual & Difference & Cluster 95\% interval \\
\midrule
Mean rank percentile & 0.4017 & 0.4296 & $+0.0279$ & $-0.0013$ to $+0.0570$ \\
Top-1 rate           & 0.0601 & 0.0622 & $+0.0021$ & $-0.0279$ to $+0.0316$ \\
Top-3 rate           & 0.1580 & 0.1118 & $-0.0462$ & $-0.0917$ to $-0.0025$ \\
Top-5 rate           & 0.2361 & 0.1641 & $-0.0720$ & $-0.1254$ to $-0.0218$ \\
\bottomrule
\end{tabular}
\end{table}

Top-3 and top-5 retrieval both fall, and both cluster intervals exclude zero. The top-5 hit rate
falls from 23.6\% to 16.4\%. Mean rank percentile moves in the unfavourable direction, from 0.4017 to
0.4296, but its cluster interval includes zero. Pair-weighted estimates point the same way:
0.4182 versus 0.4456 in mean rank percentile and 0.210 versus 0.166 in top-5 retrieval.

Both representations still beat a degree-preserving null. The null preserves every compound's
active count and every target's frequency. Absolute Vina gives mean rank percentile 0.4017 versus
0.4551 under the null and top-5 retrieval 0.2361 versus 0.1577. The two-way residual gives 0.4296
versus 0.4981 and 0.1641 versus 0.1253 ($p=0.0002$ and $p=0.0212$). Full results are in
\path{results/residual_mie_boundary_audit/null_summary.csv}. Target-specific ordering survives
removal of the shared axis, but absolute scores retrieve more actives at fixed $k$.

The algebra explains why the two-way residual cannot help here. Within one ligand,
$\Gamma_{ij} = (X_{ij} - \bar X_{\cdot j}) - c_i$, where $c_i$ is that ligand's mean of the
column-centred scores. Subtracting a constant from all of a ligand's entries cannot change their
order, so the two-way residual ranks a ligand's targets exactly as the column-centred scores do;
the maximum reconstruction error is $1.8\times10^{-15}$. It is not the same ordering as
column-$z$, which rescales each target before the subtraction: the two differ by up to 33 rank
positions on this support. A rise in spectral effective dimension after per-ligand centring can
expose covariance among targets, and it cannot create new within-ligand ranking information.

A second frozen benchmark, run on DOCKSTRING rather than Docking-44, points the same way. Over
6{,}480 ligands with 25{,}214 non-tied
target pairs on 31 targets, mean per-ligand pairwise ordering accuracy is 0.5371 for absolute
Vina, 0.5304 for column-$z$ and 0.5294 for the two-way residual; the residual-minus-absolute
difference is $-0.0077$ with a conservative 95\% interval of $-0.0220$ to $+0.0061$
(\path{results/dockstring_chembl_ranking/representation_summary.csv} and
\path{results/dockstring_chembl_ranking/contrast_summary.csv}).
